\documentclass[11pt]{article}
\usepackage{preamble}
\renewcommand{\answer}[1]{}

\begin{document}

\ladderheading{AP Statistics \textperiodcentered\ Unit 1 \textperiodcentered\ Topic 1.10 \textperiodcentered\ B: 2026-09-28 \textperiodcentered\ E: 2026-10-05}{What Can Each Study Claim?}

\focusbanner{TODAY'S PROBLEM}

Suppose a student wants to know whether a new tutoring center raises test scores.\par\smallskip \textbf{Plan A:} Compare this year's test scores for students who choose to attend the center with scores for students who do not attend.\par \textbf{Plan B:} Recruit 40 student volunteers, then use a coin flip to assign each volunteer either to attend the center or not to attend. Compare their later test scores.\par
\begin{center}
\textbf{Two plans for studying tutoring}\par\smallskip
\begin{tabular}{|l|c|c|c|}
\hline
\textbf{Plan} & \textbf{Who?} & \textbf{Groups?} & \textbf{Response} \\ \hline
\textbf{A} & Students & Choose & Scores \\ \hline
\textbf{B} & 40 volunteers & Coin flip & Scores \\ \hline
\end{tabular}
\end{center}
\begin{firsttakebox}{FIRST TAKE --- before the video (5 min)}
Which plan gives stronger evidence that the center raises scores? Name the feature behind your choice.
\workspace{0.9in}
\end{firsttakebox}

\begin{callout}{\IconBook\ SELECTION, ASSIGNMENT, AND SCOPE (keep on the page)}{calloutgreen}
An \textbf{observational study} records data without imposing treatments. An \textbf{experiment} assigns treatments. A confounding factor is an outside factor related to both the groups and the outcome. Random assignment tends to balance such factors between treatments. Random selection or other evidence of representativeness supports generalization to the population selected from. Observational comparisons alone cannot establish cause.
\end{callout}

\newpage
\textbf{Finish after the video:}\par\smallskip

\dokbadge{1}~\textbf{(a)} Identify each plan as an observational study or an experiment.\par
\workspace{0.8in}

\dokbadge{2}~\textbf{(b)} Name one outside factor that could affect both choosing tutoring and test scores in Plan A. Explain both links in two sentences.\par
\workspace{1.4in}

\dokbadge{3}~\textbf{(c)} Which plan offers stronger evidence that tutoring causes a change in scores? Explain why its design helps with the outside factor from (b). Then state why its result may not apply to all students at the school.\par
\workspace{2.0in}

\begin{sentenceframebox}\raggedright\textbf{Frame:} Plan \blank[0.4in] supports a causal claim because \blank[2.5in].\end{sentenceframebox}

\begin{sentenceframebox}\raggedright\textbf{Frame:} It still cannot generalize to \blank[1.5in] because \blank[2.2in].\end{sentenceframebox}

\begin{summaryexitbox}{EXIT --- turn this in}
One thing the video changed about my first take: \hrulefill\par
\workspace{0.4in}
\end{summaryexitbox}

\end{document}
